\begin{pblock}{Why openEO}
The archive is the bottleneck: a decade of Sentinel-2 over one estuary is
thousands of scenes, of which most are useless. Two common answers both fail
at scale — downloading whole scenes wastes almost everything fetched, and
closed cloud platforms make the processing itself unciteable.

We use \term{openEO} on the Copernicus Data Space. One batch job requests
only the bands and the polygon needed, and returns a single netCDF cube.
Everything after that is local:

\begin{itemize}
  \item \textbf{Analysis-ready L2A}, subset server-side in space, time and
        band — no local mirror of the archive.
  \item \textbf{Download once, stream forever.} Re-running an analysis never
        re-downloads; memory is bounded by the time-block size, not by the
        area or the number of years.
  \item \textbf{The cube is the unit of reproducibility}, and the process
        graph is a portable, standard document rather than a vendor's
        notebook.
\end{itemize}

\vspace{4mm}
\pfig[\linewidth]{pipeline}
\figcap{The pipeline. One cloud job, then everything local — which is what
makes it affordable to run many contiguous cells in parallel on one machine.}
\end{pblock}
